\section{Introduction}
\label{sec:intro}

% Requirement uncertainty exists and it effects generation reliability->background
Large Language Models (LLMs) have significantly advanced the state of automated code generation. Modern code generation systems are now capable of synthesizing executable programs in diverse programming languages from natural language requirements, achieving strong performance on widely adopted benchmarks such as HumanEval~\cite{humaneval}, MBPP~\cite{mbpp} and SWE-Bench~\cite{jimenez2024swebench}. Consequently, LLM-based code generation is being increasingly integrated into practical software engineering workflows through contemporary IDEs, serving a broad range of real-world development tasks across various sectors.

Despite these advances, the performance of LLM-based code generation still heavily depends on the quality of requirements. In real-world scenarios, requirements are often far from perfectly specified: they may contain vagueness, ambiguities, under-specification or implicit constraints~\cite{abdeahad2026accurate, ezzini2022anaphoric}. Such requirement-side imperfection may diverge the interpretations of code generation models, leading to requirements deviation in implementation, inconsistent outputs across different models, or instability across repeated generations, even when the underlying generation model is sufficiently capable~\cite{yang2026assessingimpactrequirementambiguity, 10554837}. In this work, we refer to these input-side deficiencies that cause unstable or low-quality code generation as \textit{requirements uncertainty}.

%Existing approaches identify it too late and too expensively,even we have clarifcation method such as clarifyGPT,it still use extra code genaration and tests...->existing method 
There is little doubt that requirement uncertainty poses a significant threat to the practical value of code intelligence.  Vijayvargiya et al.~\cite{vijayvargiya2026ambigsweinteractiveagentsovercome} reveal that AI models struggle to distinguish between well-specified and under-specified instructions, with the latter representing a common form of such uncertainty. In response, recent studies have begun to address requirement uncertainty through clarification or repair. For example, ClarifyGPT detects ambiguous requirements by generating test inputs and checking code consistency, before it asks targeted clarification questions~\cite{mu2024clarifygpt}. SpecFix repairs ambiguous problem descriptions by analyzing the distribution of programs induced by a requirement with testing and program repair techniques~\cite{jia2025specfix}. These studies provide valuable evidence that requirement uncertainty can be detected and mitigated.

However, these approaches rely on comparing consistency or distribution across multiple generation outputs, each of which is already resource-intensive. In modern code generation workflows that often involve planning, retrieval, generation, testing, iterative refinement~\cite{islam2024mapcoder,zhang2023repocoder, 11334685, zhao2025projectgen}, or multiple reasoning steps, this repeated execution cost is further multiplied. As a result, uncertainty is detected only after substantial resources have been expended across multiple runs, only to reveal that the root cause lies in the requirement itself. \emph{This motivates a lightweight pre-generation solution that assesses requirements uncertainty before any costly generation iteration is triggered.}

However, building such a solution is non-trivial.

\noindent \textbf{Challenge I: How to obtain reliable uncertainty signals for pre-generation use? }
Requirement uncertainty may arise from many forms of imperfect natural language descriptions, including ambiguity, vagueness, under-specification, implicit constraints, and even trivial noise~\cite{abdeahad2026accurate,ezzini2022anaphoric,yang2026assessingimpactrequirementambiguity}. However, not every textual imperfection necessarily harms code generation reliability. As modern LLMs become increasingly capable, some incomplete or implicit information may be inferred correctly, while seemingly minor textual variations may still lead to unstable implementations~\cite{chen2026nlperturbator}. Therefore, the key challenge is not merely to enumerate possible uncertainty types, but to obtain reliability-aware uncertainty signals that reflect whether a requirement is likely to cause unstable or low-quality code generation.

\noindent\textbf{Challenge II: How to predict requirement uncertainty before generation?} The demand of requirement uncertainty prediction before code generation makes post-generation approaches infeasible~\cite{mu2024clarifygpt, jia2025specfix}, due to the unavailability of execution results from generated code. While existing requirements ambiguity detection approaches often rely on human-labeled corpora~\cite{abdeahad2026accurate,ezzini2022anaphoric}, these methods originate from the requirements engineering community and are not tailored to code generation tasks. Our preliminary study shows that not all human-centric uncertainty uniformly transform into model-observed uncertainty.  Hence, their ambiguity labels do not directly indicate whether a requirement will lead to unreliable generation. In contrast, reliability-aware uncertainty is inherently tied to the model's generation behavior, specifically, the correctness and consistency of its outputs. Since neither the generated programs nor their execution feedback is accessible prior to generation, a practical solution must predict such uncertainty directly from the requirement text, while also implicitly taking the model's responses into account.

% \noindent\textbf{Challenge III: How to use uncertainty to guide clarification automatically and effectively?}
% Even if requirement uncertainty can be identified before generation, mitigating it automatically remains non-trivial. A natural way to reduce uncertainty is to ask targeted clarification questions and rewrite the requirement based on the received responses, as adopted by clarification-based code generation frameworks~\cite{mu2024clarifygpt}. However, effective clarification requires more than simply asking additional questions. Prior work shows that directly prompting LLMs to detect and resolve ambiguous requirements may produce irrelevant or inconsistent clarifications~\cite{jia2025specfix}. Therefore, the key challenge is how to leverage the quantified uncertainty signals to ask pertinent clarification questions and automatically rewrite the requirement based on the clarification responses, so that the clarified requirement genuinely reduces uncertainty rather than introducing redundant or misleading details.

To address the above challenges, we propose our approach \textsc{Pulse} (Pre-generation Uncertainty Labeling with Surrogate Evaluator). It includes three stages. Firstly, we inject the perturbations mentioned in literature into the requirements description of public benchmarks, and labeling the uncertainty based on the multiple execution of a moderate-size code generation model. Second, with the annotated dataset, we train a lightweight RoBERTa-based classifier as a surrogate model to directly predict uncertainty scores from surface textual features. The classifier categorizes requirements into two uncertainty levels (high- and low- risk) with low inference overhead, enabling practical pre-generation risk assessment.  The high-risk results trigger the downstream clarification process. Third, we instantiate a simple agentic workflow that adopts an asking-answering pattern to iteratively clarify high-uncertainty requirements before generation takes place.

We construct our training dataset from three standalone functional benchmarks, HumanEval-ET~\cite{humaneval}, MBPP-ET~\cite{mbpp}, and LiveCodeBench~\cite{jain2025livecodebench}, and train our classifier on it. On functional-level benchmarks with four generation approaches (two LLMs and two state-of-the-art frameworks), our experiments demonstrate: (1)  the effectiveness of \textsc{Pulse} in improving both pass@1 (by up to 46\%) and stability when requirements are clarified prior to generation; (2) the utility of our proposed uncertainty quantification (UQ) as a signal for requirement clarification, with higher UQ scores correlating with greater post-clarification gains; (3) the efficiency of our lightweight classifier. In addition, we validate the generalizability of our approach on project-level code generation tasks using a public benchmark that consists of 10 real-world repositories.

%train a lightweight classifier to predict requirement uncertainty directly from requirement text. We evaluate \textsc{Pulse} primarily on function-level code generation, and additionally conduct a preliminary project-level study to examine its applicability beyond standalone tasks. The results show that uncertainty-aware selective clarification improves pass@1 across most function-level settings and remains compatible with both direct LLM generation and representative generation-time baselines.



% latex problem -- this table should be displayed near the section Prestudy
\begin{table*}[htbp]
\centering
\caption{Requirement perturbations used in the preliminary study. These perturbations are adapted from prior studies on natural-language perturbation and prompt robustness.}
\label{tab:uncertainty_factors}
\begin{tabular*}{\textwidth}{@{\extracolsep{\fill}}p{3.3cm}p{8.0cm}p{3.0cm}p{2.0cm}@{}}
\toprule
\textbf{Perturbation Type} & \textbf{Explanation} & \textbf{Transformation Method} & \textbf{Reference} \\
\midrule

Content Rewriting
& Rewrite the requirement by abstracting or paraphrasing detailed expressions of functional constraints, edge cases, or task intent while keeping the same high-level programming task.
& LLM-assisted rewriting
& ~\cite{chen2026nlperturbator,zhu2024promptrobust} \\

Format Transformation
& Rewrite the requirement into a natural, conversational, and human-like paragraph, avoiding formal benchmark-style structures while preserving all functional requirements and constraints.
& LLM-assisted rewriting
& ~\cite{chen2026nlperturbator} \\

Example Removal
& Remove input-output examples from the requirement, such as doctest-style example usages in HumanEval and sample input/output cases in programming-contest-style tasks such as LiveCodeBench.
& Rule-based removal
& ~\cite{nashid2023retrieval,khojah2025impact} \\

Verbose Rewriting
& Add redundant descriptions or expand the requirement expression while preserving the intended functionality.
& Rule-based transformation
& ~\cite{chen2026nlperturbator,zhu2024promptrobust} \\

Noise Injection
& Insert irrelevant or distracting natural-language content into the requirement without introducing new inputs, outputs, constraints, or expected behavior.
& Rule-based transformation
& ~\cite{chen2026nlperturbator,zhu2024promptrobust} \\

Typographical Perturbation
& Introduce spelling or character-level errors while aiming to preserve the original requirement intent.
& TextAttack-based character transformation
& ~\cite{morris2020textattack,zhu2024promptrobust} \\

\bottomrule
\end{tabular*}
\end{table*}



In summary, the contributions of this paper are threefold:
\begin{itemize}[leftmargin=1em]

    \item \textbf{Approach.} We propose \textsc{Purse}, a novel framework for assessing requirement uncertainty without executing code generation. It serves as a lightweight, plug-and-play guardrail that can be seamlessly integrated with existing requirements clarification methods and code generation pipelines.
    
    \item \textbf{Labeling.} To address the scarcity of uncertainty-aware requirement datasets for reliable code generation, we propose a labeling method that annotates requirement uncertainty via execution feedback from a moderate-size code generation model, facilitating benchmark construction for both  classifier training and evaluation purposes.

    \item \textbf{Evaluation.} Extensive experiments on four public benchmarks demonstrate the effectiveness of our \textsc{Purse} and its classifier, as well as its good generalization to project-level code generation tasks. The source code and data are publicly available at  \url{https://figshare.com/s/bb1ffe8ca97f6ab5ab95}.
    % \item We introduce a model-centered definition of requirement uncertainty for code generation and characterize uncertainty through generation consistency under repeated sampling.
    % \item We propose a consistency-based uncertainty quantification method combining pass-rate entropy and CodeBLEU-based generation divergence.
    % \item We develop a lightweight requirement uncertainty classifier that enables pre-generation uncertainty prediction directly from requirement text.
    % \item We design an uncertainty-aware selective clarification framework and demonstrate its effectiveness in both function-level and repository-level code generation tasks.
\end{itemize}

